\olchapter{fol}{axd}{Axiomatic \usetoken{P}{derivation}}

\begin{editorial}
  No effort has been made yet to ensure that the material in this
  chapter respects various tags indicating which connectives and
  quantifiers are primitive or defined: all are assumed to be
  primitive, except $\liff$ which is assumed to be defined. If the FOL
  tag is true, we produce a version with quantifiers, otherwise
  without.
\end{editorial}



\olfileid{fol}{axd}{rul}

\olsection{Rules and \usetoken{P}{derivation}}

\begin{explain}
  Axiomatic !!{derivation}s are perhaps the simplest !!{derivation} system for
  logic. !!^a{derivation} is just a sequence of !!{formula}s.  To
  count as !!a{derivation}, every !!{formula} in the sequence must
  either be an instance of an axiom, or must follow from one or more
  !!{formula}s that precede it in the sequence by a rule of inference.
  !!^a{derivation} !!{derive}s its last !!{formula}.
\end{explain}

\begin{defn}[!!^{derivability}]
If $\Gamma$ is a set of !!{formula}s of $\Lang L$ then a
\emph{!!{derivation}} from~$\Gamma$ is a finite sequence $!A_1$,
\dots,~$!A_n$ of !!{formula}s where for each $i \le n$ one of the
following holds:
\begin{enumerate}
\item $!A_i \in \Gamma$; or
\item $!A_i$ is an axiom; or
\item $!A_i$ follows from some $!A_j$ (and $!A_k$) with $j < i$ (and
  $k < i$) by a rule of inference.
\end{enumerate}
\end{defn}

What counts as a correct !!{derivation} depends on which inference
rules we allow (and of course what we take to be axioms).  And an
inference rule is an if-then statement that tells us that, under
certain conditions, a step~$A_i$ in !!a{derivation} is a correct
inference step.

\begin{defn}[Rule of inference]
A \emph{rule of inference} gives a sufficient condition for what
counts as a correct inference step in !!a{derivation} from~$\Gamma$.
\end{defn}

For instance, since any one-element sequence $!A$ with $!A \in \Gamma$
trivially counts as !!a{derivation}, the following might be a very
simple rule of inference:
\begin{quote}
  If $!A \in \Gamma$, then $!A$ is always a correct inference step in
  any !!{derivation} from~$\Gamma$.
\end{quote}
Similarly, if $!A$ is one of the axioms, then $!A$ by itself is
!!a{derivation}, and so this is also a rule of inference:
\begin{quote}
  If $!A$ is an axiom, then $!A$ is a correct inference step.
\end{quote}
It gets more interesting if the rule of inference appeals to
!!{formula}s that appear before the step considered. The following
rule is called \emph{modus ponens:}
\begin{quote}
  If $!B \lif !A$ and $!B$ occur higher up in the !!{derivation},
  then~$!A$ is a correct inference step.
\end{quote}
If this is the only rule of inference, then our definition of
!!{derivation} above amounts to this: $!A_1$, \dots,~$!A_n$ is
!!a{derivation} iff for each $i \le n$ one of the following holds:
\begin{enumerate}
\item $!A_i \in \Gamma$; or
\item $!A_i$ is an axiom; or
\item for some $j < i$, $!A_j$ is $!B \lif !A_i$, and for some $k < i$,
  $!A_k$ is~$!B$.
\end{enumerate}
The last clause says that $!A_i$ follows from~$!A_j$ ($!B \lif !A_i$) and $!A_k$
($!B$) by modus ponens. If we can go from $1$ to~$n$, and
each time we find !!a{formula}~$!A_i$ that is either in~$\Gamma$, an
axiom, or which a rule of inference tells us that it is a correct
inference step, then the entire sequence counts as a correct
!!{derivation}.

\begin{defn}[!!^{derivability}]
A !!{formula}~$!A$ is \emph{!!{derivable}} from $\Gamma$, written
$\Gamma \Proves !A$, if there is !!a{derivation} from $\Gamma$ ending
in $!A$.
\end{defn}

\begin{defn}[Theorems]
!!^a{formula}~$!A$ is a \emph{theorem} if there is !!a{derivation}
of~$!A$ from the empty set.  We write $\Proves !A$ if $!A$ is a
theorem and $\Proves/ !A$ if it is not.
\end{defn}






\olfileid{fol}{axd}{prp}

\olsection{Axioms and Rules for the Propositional Connectives}

\begin{defn}[Axioms]
The set of $\PAx$ of \emph{axioms} for the propositional connectives comprises
all !!{formula}s of the following forms:
\begin{align}
  & (!A \land !B) \lif !A \ollabel{ax:land1}\\
  & (!A \land !B) \lif !B \ollabel{ax:land2}\\
  & !A \lif (!B \lif (!A \land !B)) \ollabel{ax:land3}\\
  & !A \lif (!A \lor !B) \ollabel{ax:lor1}\\
  & !A \lif (!B \lor !A) \ollabel{ax:lor2}\\
  & (!A \lif !C) \lif ((!B \lif !C) \lif ((!A \lor !B) \lif !C)) \ollabel{ax:lor3}\\
  & !A \lif (!B \lif !A) \ollabel{ax:lif1}\\
  & (!A \lif (!B \lif !C)) \lif ((!A \lif !B) \lif (!A \lif !C)) \ollabel{ax:lif2}\\
  & (!A \lif !B) \lif ((!A \lif \lnot !B) \lif \lnot !A) \ollabel{ax:lnot1}\\
  & \lnot !A \lif (!A \lif !B) \ollabel{ax:lnot2}\\
  & \ltrue \ollabel{ax:ltrue}\\
  & \lfalse \lif !A \ollabel{ax:lfalse1}\\
  & (!A \lif \lfalse) \lif \lnot !A \ollabel{ax:lfalse2}\\
  & \lnot\lnot !A \lif !A \ollabel{ax:dne}
\end{align}
\end{defn}

\begin{defn}[Modus ponens]
  If $!B$ and $!B \lif !A$ already occur in !!a{derivation}, then $!A$ is
  a correct inference step.
\end{defn}

We'll abbreviate the rule modus ponens as ``\MP.''




  

\olfileid{fol}{axd}{qua}

\olsection{Axioms and Rules for Quantifiers}

\begin{defn}[Axioms for quantifiers]
The \emph{axioms} governing quantifiers are
all instances of the following:
\begin{align}
\ollabel{ax:q1} & \lforall[x][!B] \lif !B(t), \\
\ollabel{ax:q2} & !B(t) \lif \lexists[x][!B].
\end{align}
for any closed term~$t$.
\end{defn}

\begin{defn}[Rules for quantifiers]
\item If $!B \lif !A(a)$ already occurs in the !!{derivation} and $a$
  does not occur in~$\Gamma$ or~$!B$, then $!B \lif
  \lforall[x][!A(x)]$ is a correct inference step.
\item If $!A(a) \lif !B$ already occurs in the !!{derivation} and $a$
  does not occur in~$\Gamma$ or~$!B$, then $\lexists[x][!A(x)] \lif
  !B$ is a correct inference step.
\end{defn}

We'll abbreviate either of these by ``\QR.''






\olfileid{fol}{axd}{pro}

\olsection{Examples of \usetoken{P}{derivation}}


\begin{ex}
Suppose we want to prove $(\lnot !D \lor !E) \lif (!D \lif
!E)$. Clearly, this is not an instance of any of our axioms, so we
have to use the \MP{} rule to !!{derive} it. Our only rule is~MP, which
given $!A$ and $!A \lif !B$ allows us to justify~$!B$. One
strategy would be to use \olref[prp]{ax:lor3} with $!A$ being $\lnot
!D$, $!B$ being $!E$, and $!C$ being $!D \lif !E$, i.e., the instance
\[
(\lnot !D \lif (!D \lif !E)) \lif ((!E \lif (!D \lif !E)) \lif ((\lnot
!D \lor !E) \lif (!D \lif !E))).
\]
Why? Two applications of MP yield the last part, which is what we
want.  And we easily see that $\lnot !D \lif (!D \lif !E)$ is an
instance of \olref[prp]{ax:lnot2}, and $!E \lif (!D \lif !E)$ is an
instance of \olref[prp]{ax:lif1}. So our !!{derivation} is:
\begin{derivation}
  1. &  $\lnot !D \lif (!D \lif !E)$ & \olref[prp]{ax:lnot2} \\
  2. & $(\lnot !D \lif (!D \lif !E)) \lif {}$\\
  &\qquad $((!E \lif (!D \lif !E)) \lif ((\lnot !D \lor !E) \lif (!D \lif !E)))$ & \olref[prp]{ax:lor3}\\
  3. & $(!E \lif (!D \lif !E)) \lif ((\lnot !D \lor !E) \lif (!D \lif !E))$ & 1, 2, \MP\\
  4. & $!E \lif (!D \lif !E)$ & \olref[prp]{ax:lif1}\\
  5. & $(\lnot !D \lor !E) \lif (!D \lif !E)$ & 3, 4, \MP
\end{derivation}
\end{ex}

\begin{ex}\ollabel{ex:identity}
Let's try to find !!a{derivation} of $!D \lif !D$.  It is not an
instance of an axiom, so we have to use \MP{} to !!{derive} it.
\olref[prp]{ax:lif1} is an axiom of the form~$!A \lif !B$ to which we
could apply~\MP. To be useful, of course, the $!B$ which \MP{} would
justify as a correct step in this case would have to be~$!D \lif !D$,
since this is what we want to !!{derive}. That means $!A$ would also
have to be $!D$, i.e., we might look at this instance of
\olref[prp]{ax:lif1}:
\[
!D \lif (!D \lif !D)
\]
In order to apply \MP, we would also need to justify the corresponding
second premise, namely~$!A$. But in our case, that would be~$!D$, and
we won't be able to !!{derive}~$!D$ by itself. So we need a different
strategy.

The other axiom involving just~$\lif$ is \olref[prp]{ax:lif2}, i.e.,
\[
(!A \lif (!B \lif !C)) \lif ((!A \lif !B) \lif (!A \lif !C))
\]
We could get to the last nested conditional by applying \MP{}
twice. Again, that would mean that we want an instance of
\olref[prp]{ax:lif2} where $!A \lif !C$ is $!D \lif !D$, the !!{formula} we
are aiming for. Then of course, $!A$ and $!C$ are both~$!D$. How
should we pick~$!B$ so that both $!A \lif (!B \lif !C)$ and $!A \lif
!B$, i.e., in our case $!D \lif (!B \lif !D)$ and $!D \lif !B$, are
also !!{derivable}? Well, the first of these is already an instance of
\olref[prp]{ax:lif1}, whatever we decide $!B$ to be. And $!D \lif !B$ would
be another instance of \olref[prp]{ax:lif1} if $!B$ were $(!D \lif !D)$.
So, our !!{derivation} is:
\begin{derivation}
  1. & $!D \lif ((!D \lif !D) \lif !D)$ & \olref[prp]{ax:lif1}\\
  2. & $(!D \lif ((!D \lif !D) \lif !D)) \lif {}$\\
  & \qquad $((!D \lif (!D \lif !D)) \lif (!D \lif !D))$ & \olref[prp]{ax:lif2}\\
  3. & $(!D \lif (!D \lif !D)) \lif (!D \lif !D)$ & 1, 2, \MP\\
  4. & $!D \lif (!D \lif !D)$ & \olref[prp]{ax:lif1}\\
  5. & $!D \lif !D$ & 3, 4, \MP
\end{derivation}
\end{ex}

\begin{ex}\ollabel{ex:chain}
Sometimes we want to show that there is !!a{derivation} of some
!!{formula} from some other !!{formula}s~$\Gamma$. For instance, let's
show that we can !!{derive} $!A \lif !C$ from $\Gamma = \{!A \lif !B,
!B \lif !C\}$.
\begin{derivation}
  1. & $!A \lif !B$ & \Hyp\\
  2. & $!B \lif !C$ & \Hyp\\
  3. & $(!B \lif !C) \lif (!A \lif (!B \lif !C))$ & \olref[prp]{ax:lif1} \\
  4. & $!A \lif (!B \lif !C)$ & 2, 3, \MP\\
  5. & $(!A \lif (!B \lif !C)) \lif {}$\\
  & \qquad $((!A \lif !B) \lif (!A \lif !C))$ & \olref[prp]{ax:lif2}\\
  6. & $((!A \lif !B) \lif (!A \lif !C))$ & 4, 5, \MP\\
  7. & $!A \lif !C$ & 1, 6, \MP
\end{derivation}
The lines labelled ``\Hyp'' (for ``hypothesis'') indicate that the
!!{formula} on that line is !!a{element} of~$\Gamma$.
\end{ex}

\begin{prop}
  \ollabel{prop:chain} If $\Gamma \Proves !A \lif !B$ and $\Gamma
  \Proves !B \lif !C$, then $\Gamma \Proves !A \lif !C$
\end{prop}

\begin{proof}
  Suppose $\Gamma \Proves !A \lif !B$ and $\Gamma \Proves !B \lif
  !C$. Then there is !!a{derivation} of $!A \lif !B$ from~$\Gamma$;
  and !!a{derivation} of~$!B \lif !C$ from~$\Gamma$ as well. Combine
  these into a single !!{derivation} by concatenating them. Now add
  lines 3--7 of the !!{derivation} in the preceding example. This is
  !!a{derivation} of $!A \lif !C$---which is the last line of the new
  !!{derivation}---from~$\Gamma$. Note that the justifications of
  lines 4 and~7 remain valid if the reference to line number~1 is
  replaced by reference to the last line of the !!{derivation} of~$!A
  \lif !B$, and reference to line number~2 by reference to the last
  line of the !!{derivation} of~$!B \lif !C$.
\end{proof}

\begin{prob} 
Show that the following hold by exhibiting !!{derivation}s from the
axioms:
\begin{enumerate} 
\item $(!A \land !B) \lif (!B \land !A)$
\item $((!A \land !B) \lif !C) \lif (!A \lif (!B \lif !C))$
\item $\lnot(!A \lor !B) \lif \lnot !A$
\end{enumerate} 
\end{prob}






  

\olfileid{fol}{axd}{prq}

\olsection{\usetoken{P}{derivation} with Quantifiers}

\begin{ex}
Let us give !!a{derivation} of $(\lforall[x][!A(x)] \land
\lforall[y][!B(y)]) \lif \lforall[x][(!A(x) \land !B(x))]$.

First, note that
\begin{align*}
  (\lforall[x][!A(x)] \land \lforall[y][!B(y)]) & \lif \lforall[x][!A(x)]\\
  \intertext{is an instance of \olref[prp]{ax:land1}, and}
  \lforall[x][!A(x)] & \lif !A(a) \\
  \intertext{of \olref[qua]{ax:q1}. So, by \olref[pro]{prop:chain}, we know that}
  (\lforall[x][!A(x)] \land \lforall[y][!B(y)]) & \lif !A(a) \\
  \intertext{is !!{derivable}. Likewise, since}
  (\lforall[x][!A(x)] \land \lforall[y][!B(y)]) & \lif \lforall[y][!B(y)] \qquad\text{and}\\
    \lforall[y][!B(y)] & \lif !B(a)\\
    \intertext{are instances of \olref[prp]{ax:land2} and \olref[qua]{ax:q1}, respectively,}
    (\lforall[x][!A(x)] \land \lforall[y][!B(y)]) & \lif !B(a)\\
    \intertext{is derivable by \olref[pro]{prop:chain}. Using an appropriate instance of \olref[prp]{ax:land3} and two applications of~\MP, we see that}
    (\lforall[x][!A(x)] \land \lforall[y][!B(y)]) & \lif (!A(a) \land !B(a))\\
    \intertext{is derivable. We can now apply \QR{} to obtain}
(\lforall[x][!A(x)] \land \lforall[y][!B(y)]) & \lif \lforall[x][(!A(x) \land !B(x))].
\end{align*}
\end{ex}






\olfileid{fol}{axd}{ptn}

\olsection{Proof-Theoretic Notions}

\begin{explain}
Just as we've defined a number of important semantic notions
(validity, entailment, satisfiability), we now
define corresponding \emph{proof-theoretic notions}.  These are not
defined by appeal to satisfaction of !!{sentence}s in !!{structure}s,
but by appeal to the !!{derivability} or !!{nonderivability} of
certain formulas.  It was an important discovery that these notions
coincide.  That they do is the content of the \emph{soundness} and
\emph{completeness theorems}.
\end{explain}

\begin{defn}[!!^{derivability}]
!!^a{formula}~$!A$ is \emph{!!{derivable}} from $\Gamma$, written
$\Gamma \Proves !A$, if there is !!a{derivation} from~$\Gamma$ ending
in~$!A$.
\end{defn}

\begin{defn}[Theorems]
!!^a{formula}~$!A$ is a \emph{theorem} if there is !!a{derivation} of
$!A$ from the empty set.  We write $\Proves !A$ if $!A$ is a theorem
and $\Proves/ !A$ if it is not.
\end{defn}

\begin{defn}[Consistency]
A set $\Gamma$ of !!{formula}s is \emph{consistent} if and only if
$\Gamma\Proves/ \lfalse$; it is \emph{inconsistent} otherwise.
\end{defn}

\begin{prop}[Reflexivity]
\ollabel{prop:reflexivity}
If $!A \in \Gamma$, then $\Gamma \Proves !A$.
\end{prop}

\begin{proof}
  The !!{formula}~$!A$ by itself is !!a{derivation} of~$!A$ from~$\Gamma$.
\end{proof}

\begin{prop}[Monotonicity]
\ollabel{prop:monotonicity}
If $\Gamma \subseteq \Delta$ and $\Gamma \Proves !A$, then $\Delta
\Proves !A$.
\end{prop}

\begin{proof}
Any !!{derivation} of $!A$ from $\Gamma$ is also !!a{derivation} of
$!A$ from~$\Delta$.
\end{proof}

\begin{prop}[Transitivity]
\ollabel{prop:transitivity}
If $\Gamma \Proves !A$ and $\{!A\} \cup \Delta \Proves
!B$, then $\Gamma \cup \Delta \Proves !B$.
\end{prop}

\begin{proof}
  Suppose $\{!A\} \cup \Delta \Proves !B$. Then there is
  !!a{derivation} $!B_1$, \dots, $!B_l = !B$ from~$\{!A\} \cup
  \Delta$. Some of the steps in that !!{derivation} will be correct
  because of a rule which refers to a prior line~$!B_i = !A$. By
  hypothesis, there is !!a{derivation} of~$!A$ from~$\Gamma$, i.e.,
  !!a{derivation}~$!A_1$, \dots, $!A_k = !A$ where every $!A_i$ is an
  axiom, !!a{element} of~$\Gamma$, or correct by a rule of
  inference. Now consider the sequence
  \[
  !A_1, \dots, !A_k = !A, !B_1, \dots, !B_l = !B.
  \]
  This is a correct !!{derivation} of~$!B$ from $\Gamma \cup \Delta$
  since every $B_i = !A$ is now justified by the same rule which
  justifies~$!A_k = !A$.
\end{proof}

Note that this means that in particular if $\Gamma \Proves !A$ and $!A
\Proves !B$, then $\Gamma \Proves !B$. It follows also that if $!A_1,
\dots, !A_n \Proves !B$ and $\Gamma \Proves !A_i$ for each~$i$, then
$\Gamma \Proves !B$.

\begin{prop}
\ollabel{prop:incons}
$\Gamma$ is inconsistent iff $\Gamma \Proves !A$ for every~$!A$.
\end{prop}

\begin{proof}
Exercise.
\end{proof}


\begin{prob}
Prove \olref[fol][axd][ptn]{prop:incons}.
\end{prob}




\begin{prop}[Compactness]
\ollabel{prop:proves-compact}
  \begin{enumerate}
  \item If $\Gamma \Proves !A$ then there is a finite subset $\Gamma_0
    \subseteq \Gamma$ such that $\Gamma_0 \Proves !A$.
  \item If every finite subset of~$\Gamma$ is
    consistent, then $\Gamma$ is consistent.
  \end{enumerate}
\end{prop}

\begin{proof}
  \begin{enumerate}
    \item If $\Gamma \Proves !A$, then there is a finite sequence of
      !!{formula}s $!A_1$, \dots,~$!A_n$ so that $!A \ident !A_n$ and
      each $!A_i$ is either a logical axiom, !!a{element} of~$\Gamma$
      or follows from previous !!{formula}s by modus ponens.  Take
      $\Gamma_0$ to be those $!A_i$ which are in~$\Gamma$.  Then the
      !!{derivation} is likewise !!a{derivation} from~$\Gamma_0$, and
      so $\Gamma_0 \Proves !A$.
    \item This is the contrapositive of~(1) for the special case $!A
      \ident \lfalse$.
\end{enumerate}
\end{proof}





\olfileid{fol}{axd}{ded}

\olsection{The Deduction Theorem}

As we've seen, giving !!{derivation}s in an axiomatic system is
cumbersome, and !!{derivation}s may be hard to find. Rather than
actually write out long lists of !!{formula}s, it is generally easier
to argue that such !!{derivation}s exist, by making use of a few
simple results. We've already established three such results:
\olref[ptn]{prop:reflexivity} says we can always assert that $\Gamma
\Proves !A$ when we know that $!A \in
\Gamma$. \olref[ptn]{prop:monotonicity} says that if $\Gamma \Proves !A$
then also $\Gamma \cup \{!B\} \Proves !A$. And
\olref[ptn]{prop:transitivity} implies that if $\Gamma \Proves !A$ and
$!A \Proves !B$, then $\Gamma \Proves !B$. Here's another simple
result, a ``meta''-version of modus ponens:

\begin{prop}
  \ollabel{prop:mp} If $\Gamma \Proves !A$ and $\Gamma \Proves !A \lif
  !B$, then $\Gamma \Proves !B$.
\end{prop}

\begin{proof}
  We have that $\{!A, !A \lif !B\} \Proves !B$:
  \begin{derivation}
    1. & $!A$ & Hyp.\\
    2. & $!A \lif !B$ & Hyp.\\
    3. & $!B$ & 1, 2, MP
  \end{derivation}
  By \olref[ptn]{prop:transitivity}, $\Gamma \Proves !B$.
\end{proof}

The most important result we'll use in this context is the deduction
theorem:

\begin{thm}[Deduction Theorem]
\ollabel{thm:deduction-thm} $\Gamma \cup \{!A\} \Proves !B$ if and
only if $\Gamma \Proves !A \lif !B$.
\end{thm}

\begin{proof}
The ``if'' direction is immediate.  If $\Gamma \Proves !A \lif !B$
then also $\Gamma \cup \{!A\} \Proves !A \lif !B$ by
\olref[ptn]{prop:monotonicity}.  Also, $\Gamma \cup \{!A\} \Proves !A$ by
\olref[ptn]{prop:reflexivity}. So, by \olref{prop:mp}, $\Gamma \cup
\{!A\} \Proves !B$.

For the ``only if'' direction, we proceed by induction on the length
of the !!{derivation} of $!B$ from $\Gamma \cup \{!A\}$.

For the induction basis, we prove the claim for every !!{derivation}
of length~$1$. !!^a{derivation} of~$!B$ from $\Gamma \cup \{!A\}$ of
length~$1$ consists of $!B$ by itself; and if it is correct $!B$ is
either $\in \Gamma \cup \{!A\}$ or is an axiom.  If $!B \in \Gamma$ or
is an axiom, then $\Gamma \Proves !B$. We also have that $\Gamma
\Proves !B \lif (!A \lif !B)$ by \olref[prp]{ax:lif1}, and
\olref{prop:mp} gives $\Gamma \Proves !A \lif !B$. If $!B \in \{ !A\}$
then $\Gamma \Proves !A \lif !B$ because the last !!{sentence}~$!A
\lif !B$ is the same as $!A \lif !A$, and we have !!{derive}d that in
\olref[pro]{ex:identity}.

For the inductive step, suppose !!a{derivation} of~$!B$ from
$\Gamma \cup \{!A\}$ ends with a step~$!B$ which is justified by modus
ponens. (If it is not justified by modus ponens, $!B \in \Gamma$, $!B
\ident !A$, or $!B$ is an axiom, and the same reasoning as in the
induction basis applies.) Then some previous steps in the
!!{derivation} are $!C \lif !B$ and $!C$, for some !!{formula}~$!C$,
i.e., $\Gamma \cup \{!A\} \Proves !C \lif !B$ and $\Gamma \cup \{!A\}
\Proves !C$, and the respective !!{derivation}s are shorter, so the
inductive hypothesis applies to them. We thus have both:
\begin{align*}
  & \Gamma \Proves !A \lif (!C \lif !B); \\
  & \Gamma \Proves !A \lif !C.
\end{align*}
But also
\[
\Gamma \Proves (!A \lif (!C \lif !B)) \lif
((!A\lif !C)  \lif (!A \lif !B)),
\]
by \olref[prp]{ax:lif2}, and two applications of \olref{prop:mp} give
$\Gamma \Proves !A \lif !B$, as required.
\end{proof}

Notice how \olref[prp]{ax:lif1} and \olref[prp]{ax:lif2} were chosen
precisely so that the Deduction Theorem would hold.

The following are some useful facts about !!{derivability}, which we
leave as exercises.

\begin{prop}
  \ollabel{prop:derivfacts}
\begin{enumerate}
\item $\Proves (!A \lif !B) \lif ((!B \lif !C)
  \lif (!A \lif !C)$; \ollabel{derivfacts:a}
\item If $\Gamma \cup \{ \lnot !A\}
  \Proves \lnot !B$ then $\Gamma \cup \{ !B\} \Proves
  !A$ (Contraposition); \ollabel{derivfacts:b}
\item  $\{ !A, \lnot!A\} \Proves
    !B$ (Ex Falso Quodlibet, Explosion); \ollabel{derivfacts:c}
\item  $\{ \lnot\lnot!A\} \Proves
  !A$ (Double Negation Elimination);\ollabel{derivfacts:d}
\item If $\Gamma \Proves \lnot\lnot!A$ then $\Gamma \Proves
  !A$;\ollabel{derivfacts:e}
\end{enumerate}
\end{prop}


\begin{prob}
  Prove \olref[fol][axd][ded]{prop:derivfacts}
\end{prob}







  

\olfileid{fol}{axd}{ddq}

\olsection{The Deduction Theorem with Quantifiers}

\begin{thm}[Deduction Theorem]
\ollabel{thm:deduction-thm-q} If $\Gamma \cup \{!A\} \Proves !B$,
then $\Gamma \Proves !A \lif !B$.
\end{thm}

\begin{proof}
We again proceed by induction on the length
of the !!{derivation} of $!B$ from $\Gamma \cup \{!A\}$.

The proof of the induction basis is identical to that in the proof
of~\olref[ded]{thm:deduction-thm}.

For the inductive step, suppose again that the !!{derivation} of $!B$
from $\Gamma \cup \{!A\}$ ends with a step~$!B$ which is justified by
an inference rule. If the inference rule is modus ponens, we proceed
as in the proof of \olref[ded]{thm:deduction-thm}. If the inference
rule is \QR, we know that $!B \ident !C \lif \lforall[x][!D(x)]$ and
!!a{formula} of the form $!C \lif !D(a)$ appears earlier in the
!!{derivation}, where $a$ does not occur in~$!C$, $!A$, or $\Gamma$. We
thus have that
\begin{align*}
  \Gamma \cup \{!A\} & \Proves !C \lif !D(a),\\
  \intertext{and the induction hypothesis applies, i.e., we have that}
    \Gamma & \Proves !A \lif (!C \lif !D(a)).\\
  \intertext{By}
  & \Proves (!A \lif (!C \lif !D(a))) \lif ((!A \land !C) \lif !D(a))\\
  \intertext{and modus ponens we get}
  \Gamma & \Proves (!A \land !C) \lif !D(a).\\
  \intertext{Since the eigenvariable condition still applies, we can add a step to this !!{derivation} justified by \QR, and get}
    \Gamma & \Proves (!A \land !C) \lif \lforall[x][!D(x)].\\
    \intertext{We also have}
    & \Proves ((!A \land !C) \lif \lforall[x][!D(x)]) \lif (!A \lif (!C \lif \lforall[x][!D(x)]),\\
    \intertext{so by modus ponens,}
    \Gamma & \Proves !A \lif (!C \lif \lforall[x][!D(x)]),
\end{align*}
i.e., $\Gamma \Proves !B$.

We leave the case where $!B$ is justified by the rule \QR, but is of
the form $\lexists[x][!D(x)] \lif !C$, as an exercise.
\end{proof}

\begin{prob}
  Complete the proof of \olref[fol][axd][ddq]{thm:deduction-thm-q}.
\end{prob}






\olfileid{fol}{axd}{prv}

\olsection{\usetoken{S}{derivability} and Consistency}

We will now establish a number of properties of the !!{derivability}
relation.  They are independently interesting, but each will play a
role in the proof of the completeness theorem.

\begin{prop}\ollabel{prop:provability-contr}
  If $\Gamma \Proves !A$ and $\Gamma \cup \{!A\}$ is inconsistent,
  then $\Gamma$ is inconsistent.
\end{prop}

\begin{proof}
  If $\Gamma \cup \{!A\}$ is inconsistent, then $\Gamma \cup \{!A\}
  \Proves \lfalse$. By \olref[ptn]{prop:reflexivity}, $\Gamma \Proves
  !B$ for every $!B \in \Gamma$. Since also $\Gamma \Proves !A$ by
  hypothesis, $\Gamma \Proves !B$ for every $!B \in \Gamma \cup
  \{!A\}$. By \olref[ptn]{prop:transitivity}, $\Gamma \Proves
  \lfalse$, i.e., $\Gamma$ is inconsistent.
\end{proof}

\begin{prop}
\ollabel{prop:prov-incons}
$\Gamma \Proves !A$ iff $\Gamma \cup \{\lnot !A\}$ is inconsistent.
\end{prop}

\begin{proof}
First suppose $\Gamma \Proves !A$. Then $\Gamma \cup \{\lnot !A\}
\Proves !A$ by \olref[ptn]{prop:monotonicity}. $\Gamma \cup \{\lnot
!A\} \Proves \lnot !A$ by \olref[ptn]{prop:reflexivity}.  We also have
$\Proves \lnot !A \lif (!A \lif \lfalse)$ by \olref[prp]{ax:lnot2}. So
by two applications of \olref[ded]{prop:mp}, we have $\Gamma \cup \{\lnot
!A\} \Proves \lfalse$.

Now assume $\Gamma \cup \{\lnot !A\}$ is inconsistent, i.e., $\Gamma
\cup \{\lnot !A\} \Proves \lfalse$. By the deduction theorem, $\Gamma
\Proves \lnot !A \lif \lfalse$. $\Gamma \Proves (\lnot !A \lif
\lfalse) \lif \lnot\lnot !A$ by \olref[prp]{ax:lfalse2}, so $\Gamma
\Proves \lnot\lnot !A$ by \olref[ded]{prop:mp}. Since $\Gamma \Proves
\lnot\lnot !A \lif !A$ (\olref[prp]{ax:dne}), we have $\Gamma
\Proves !A$ by \olref[ded]{prop:mp} again.
\end{proof}

\begin{prob}
Prove that $\Gamma \Proves \lnot !A$ iff $\Gamma \cup \{!A\}$ is
inconsistent.
\end{prob}

\begin{prop}\ollabel{prop:explicit-inc}
  If $\Gamma \Proves !A$ and $\lnot !A \in \Gamma$, then $\Gamma$ is
  inconsistent.
\end{prop}

\begin{proof}
  $\Gamma \Proves \lnot !A \lif (!A \lif \lfalse)$ by \olref[prp]{ax:lnot2}.
  $\Gamma \Proves \lfalse$ by two applications of \olref[ded]{prop:mp}.
\end{proof}


\begin{prop}\ollabel{prop:provability-exhaustive}
  If $\Gamma \cup \{!A\}$ and $\Gamma \cup \{\lnot !A\}$ are both
  inconsistent, then $\Gamma$ is inconsistent.
\end{prop}

\begin{proof}
Exercise.
\end{proof}


\begin{prob}
  Prove \olref[fol][axd][prv]{prop:provability-exhaustive}
\end{prob}









\olfileid{fol}{axd}{ppr}

\olsection{\usetoken{S}{derivability} and the Propositional Connectives}

\begin{explain}
  We establish that the !!{derivability} relation~$\Proves$ of
  axiomatic deduction is strong enough to establish some basic facts
  involving the propositional connectives, such as that $!A \land !B
  \Proves !A$ and $!A, !A \lif !B \Proves !B$ (modus ponens). These
  facts are needed for the proof of the completeness theorem.
\end{explain}

\begin{prop}\ollabel{prop:provability-land}
  \begin{enumerate}
  \item \ollabel{prop:provability-land-left} Both $!A \land !B \Proves
    !A$ and $!A \land !B \Proves !B$
  \item \ollabel{prop:provability-land-right} $!A, !B \Proves !A \land !B$.
  \end{enumerate}
\end{prop}

\begin{proof}
  \begin{enumerate}
    \item From \olref[prp]{ax:land1} and \olref[prp]{ax:land1} by
      modus ponens.
  \item From \olref[prp]{ax:land3} by two applications of
    modus ponens.
  \end{enumerate}
\end{proof}


\begin{prop}\ollabel{prop:provability-lor}
  \begin{enumerate}
  \item $!A \lor !B, \lnot !A, \lnot !B$ is inconsistent.
  \item Both $!A \Proves !A \lor !B$ and $!B \Proves !A \lor !B$.
  \end{enumerate}
\end{prop}

\begin{proof}
  \begin{enumerate}
  \item From \olref[prp]{ax:lnot1} we get $\Proves \lnot !A \lif (!A
    \lif \lfalse)$ and $\Proves \lnot !B \lif (!B \lif \lfalse)$. So
    by the deduction theorem, we have $\{\lnot !A\} \Proves !A \lif
    \lfalse$ and $\{\lnot !B\} \Proves !B \lif \lfalse$. From
    \olref[prp]{ax:lor3} we get $\{\lnot !A, \lnot !B\} \Proves (!A
    \lor !B) \lif \lfalse$. By the deduction theorem, $\{!A \lor !B,
    \lnot !A, \lnot !B\} \Proves \lfalse$.
  \item From \olref[prp]{ax:lor1} and \olref[prp]{ax:lor2} by modus
    ponsens.
  \end{enumerate}
\end{proof}

\begin{prop}\ollabel{prop:provability-lif}
  \begin{enumerate}
  \item \ollabel{prop:provability-lif-left}  $!A, !A \lif !B \Proves !B$.
  \item \ollabel{prop:provability-lif-right}
    Both $\lnot !A \Proves !A \lif !B$ and $!B \Proves !A \lif !B$.
  \end{enumerate}
\end{prop}

\begin{proof}
  \begin{enumerate}
  \item We can !!{derive}:
    \begin{derivation}
      1. & $!A$ & \Hyp\\
      2. & $!A \lif !B$ & \Hyp\\
      3. & $!B$ & 1, 2, \MP
    \end{derivation}
    
  \item By \olref[prp]{ax:lnot2} and \olref[prp]{ax:lif1} and the
    deduction theorem, respectively.
  \end{enumerate}
\end{proof}




  

\olfileid{fol}{axd}{qpr}

\olsection{\usetoken{S}{derivability} and the Quantifiers}

\begin{explain}
  The completeness theorem also requires that axiomatic deductions
  yield the facts about~$\Proves$ established in this section.
\end{explain}

\begin{thm}
\ollabel{thm:strong-generalization} If $c$ is !!a{constant} not occurring
in $\Gamma$ or $!A(x)$ and $\Gamma \Proves !A(c)$, then $\Gamma
\Proves \lforall[x][!A(x)]$.
\end{thm}

\begin{proof}
By the deduction theorem, $\Gamma \Proves \ltrue \lif !A(c)$. Since
$c$ does not occur in $\Gamma$ or~$\top$, we get $\Gamma \Proves
\ltrue \lif \lforall[x][!A(x)]$. By the deduction theorem again, $\Gamma \Proves
\lforall[x][!A(x)]$.
\end{proof}

\begin{prop}
\ollabel{prop:provability-quantifiers}
\begin{tagenumerate}{prvEx,prvAll}
$!A(t) \Proves \lexists[x][!A(x)]$.

$\lforall[x][!A(x)] \Proves !A(t)$.
\end{tagenumerate}
\end{prop}

\begin{proof}
\begin{tagenumerate}{prvEx,prvAll}
By \olref[qua]{ax:q2} and the deduction theorem.
By \olref[qua]{ax:q1} and the deduction theorem.
\end{tagenumerate}
\end{proof}






\olfileid{fol}{axd}{sou}
      
\olsection{Soundness}

\begin{explain}
!!^a{derivation} system, such as axiomatic deduction, is \emph{sound}
if it cannot !!{derive} things that do not actually hold.  Soundness is
thus a kind of guaranteed safety property for !!{derivation} systems.
Depending on which proof theoretic property is in question, we would
like to know for instance, that
\begin{enumerate}
\item every !!{derivable}~$!A$ is valid;
\item if $!A$ is !!{derivable} from some others~$\Gamma$, it is also a
  consequence of them;
\item if a set of !!{formula}s~$\Gamma$ is inconsistent, it is
  unsatisfiable.
\end{enumerate}
These are important properties of !!a{derivation} system.  If any of them do
not hold, the !!{derivation} system is deficient---it would !!{derive} too much.
Consequently, establishing the soundness of !!a{derivation} system is of the
utmost importance.
\end{explain}

\begin{prop}
  If $!A$ is an axiom, then
  $\Sat{M}{!A}[s]$ for each !!{structure}~$\Struct{M}$ and assignment~$s$.
\end{prop}

\begin{proof}
We have to verify that all the axioms are valid. For
  instance, here is the case for \olref[qua]{ax:q1}: suppose $t$ is
  !!{free for} $x$ in $!A$, and assume
  $\Sat{M}{\lforall[x][!A]}[s]$. Then by definition of satisfaction,
  for each $\varAssign{s'}{s}{x}$, also $\Sat{M}{!A}[s']$, and in particular
  this holds when $s'(x) = \Value{t}{M}[s]$. By
  \olref[syn][ext]{prop:ext-formulas},
  $\Sat{M}{\Subst{!A}{t}{x}}[s]$. This shows that
  $\Sat{M}{(\lforall[x][!A] \lif \Subst{!A}{t}{x})}[s]$.
\end{proof}

\begin{thm}[Soundness]
\ollabel{thm:soundness}
If $\Gamma \Proves !A$ then $\Gamma \Entails !A$.
\end{thm}

\begin{proof}
By induction on the length of the !!{derivation} of $!A$ from
$\Gamma$. If there are no steps justified by inferences, then all
!!{formula}s in the !!{derivation} are either instances of axioms or are
in~$\Gamma$. By the previous proposition, all the axioms are
valid, and hence if $!A$ is an axiom then
$\Gamma \Entails !A$. If $!A \in \Gamma$, then trivially $\Gamma
\Entails !A$.

If the last step of the !!{derivation} of~$!A$ is justified by modus
ponens, then there are !!{formula}s $!B$ and $!B \lif !A$ in the
!!{derivation}, and the induction hypothesis applies to the part of
the !!{derivation} ending in those !!{formula}s (since they contain at
least one fewer step justified by an inference).  So, by induction
hypothesis, $\Gamma \Entails !B$ and $\Gamma \Entails !B \lif
!A$. Then $\Gamma \Entails !A$ by
\olref[syn][sem]{thm:sem-deduction}.

Now suppose the last step is justified by
\QR. Then that step has the form $!C \lif \lforall[x][B(x)]$ and
there is a preceding step $!C \lif !B(c)$ with $c$ not in $\Gamma$,
$!C$, or $\lforall[x][B(x)]$. By induction hypothesis, $\Gamma
\Entails !C \lif !B(c)$. By
\olref[syn][sem]{thm:sem-deduction}, $\Gamma \cup \{!C\} \Entails
!B(c)$.

Consider some structure~$\Struct{M}$ such that $\Sat{M}{\Gamma \cup
  \{!C\}}$.  We need to show that $\Sat{M}{\lforall[x][!B(x)]}$. Since
$\lforall[x][!B(x)]$ is !!a{sentence}, this means we have to show that
for every variable assignment~$s$, $\Sat{M}{!B(x)}[s]$
(\olref[syn][ass]{prop:sat-quant}). Since $\Gamma \cup \{!C\}$
consists entirely of sentences, $\Sat{M}{!D}[s]$ for all $!D \in
\Gamma$ by \olref[syn][sat]{defn:satisfaction}.  Let $\Struct{M'}$ be
like $\Struct{M}$ except that $\Assign{c}{M'} = s(x)$.  Since $c$ does
not occur in~$\Gamma$ or~$!C$, $\Sat{M'}{\Gamma \cup \{!C\}}$ by
\olref[syn][ext]{cor:extensionality-sent}. Since $\Gamma \cup \{!C\}
\Entails !B(c)$, $\Sat{M'}{B(c)}$.  Since $!B(c)$ is !!a{sentence},
$\Sat{M'}{!B(c)}[s]$ by
\olref[syn][ass]{prop:sentence-sat-true}. $\Sat{M'}{!B(x)}[s]$ iff
$\Sat{M'}{!B(c)}[s]$ by \olref[syn][ext]{prop:ext-formulas} (recall that
$!B(c)$ is just $\Subst{!B(x)}{c}{x}$). So,
$\Sat{M'}{!B(x)}[s]$. Since $c$ does not occur in~$!B(x)$, by
\olref[syn][ext]{prop:extensionality}, $\Sat{M}{!B(x)}[s]$. But $s$
was an arbitrary variable assignment, so
$\Sat{M}{\lforall[x][!B(x)]}$. Thus $\Gamma \cup \{!C\} \Entails
\lforall[x][!B(x)]$. By \olref[syn][sem]{thm:sem-deduction}, $\Gamma
\Entails !C \lif \lforall[x][!B(x)]$.

The case where $!A$ is justified by \QR{} but is of the form
$\lexists[x][!B(x)] \lif !C$ is left as an exercise.
\end{proof}


\begin{prob}
  Complete the proof of \olref[fol][axd][sou]{thm:soundness}.
\end{prob}


\begin{cor}
\ollabel{cor:weak-soundness}
If $\Proves !A$, then $!A$ is valid.
\end{cor}

\begin{cor}
\ollabel{cor:consistency-soundness}
If $\Gamma$ is satisfiable, then it is consistent.
\end{cor}

\begin{proof}
We prove the contrapositive.  Suppose that $\Gamma$ is not consistent.
Then $\Gamma \Proves \lfalse$, i.e., there is !!a{derivation} of
$\lfalse$ from~$\Gamma$. By \olref{thm:soundness}, any
!!{structure}~$\Struct{M}$
that satisfies $\Gamma$ must satisfy~$\lfalse$.  Since
$\Sat/{M}{\lfalse}$ for every
  !!{structure}~$\Struct{M}$, no
$\Struct{M}$ can satisfy $\Gamma$, i.e.,
$\Gamma$ is not satisfiable.
\end{proof}





  

\olfileid{fol}{axd}{ide}

\olsection{\usetoken{P}{derivation} with \usetoken{S}{identity}}

In order to accommodate $\eq$ in !!{derivation}s, we simply add new
axiom schemas. The definition of !!{derivation} and $\Proves$ remains
the same, we just also allow the new axioms.

\begin{defn}[Axioms for !!{identity}]
\begin{align}
\ollabel{ax:id1} & \eq[t][t], \\
\ollabel{ax:id2} & \eq[t_1][t_2] \lif (!B(t_1) \lif
  !B(t_2)),
\end{align}
for any closed terms $t$, $t_1$, $t_2$.
\end{defn}

\begin{prop}\ollabel{prop:sound}
  The axioms \olref{ax:id1} and \olref{ax:id2} are valid.
\end{prop}

\begin{proof}
  Exercise.
\end{proof}

\begin{prob}
Prove \olref[fol][axd][ide]{prop:sound}.
\end{prob}

\begin{prop}\ollabel{prop:iden1}
 $\Gamma \Proves \eq[t][t]$, for any term $t$ and set~$\Gamma$.
\end{prop}

\begin{prop}\ollabel{prop:iden2}
  If $\Gamma \Proves !A(t_1)$ and $\Gamma \Proves
  \eq[t_1][t_2]$, then $\Gamma \Proves !A(t_2)$.
\end{prop}

\begin{proof}
The !!{formula}
\[
(\eq[t_1][t_2] \lif (!A(t_1) \lif !A(t_2)))
\]
is an instance of~\olref{ax:id2}. The conclusion follows by two applications
of~\MP.
\end{proof}




\OLEndChapterHook






